\documentclass[10pt,a4paper]{amsart}
\usepackage[margin=19mm]{geometry}
\usepackage{amsmath,amssymb,mathtools}
\usepackage[scr=rsfs,cal=euler]{mathalfa}
\usepackage{xfrac}
\usepackage[unicode,hidelinks,bookmarks=false]{hyperref}
\newcommand{\ie}{i.e.\ }
\newcommand{\cf}{cf.\ }
\newcommand{\resp}{respectively\ }
\newcommand{\Subsection}{\subsection}
\providecommand{\nobreakdash}{\nobreak\hbox{-}}
\setlength{\emergencystretch}{2em}
\multlinegap0pt
\usepackage{siunitx}
\usepackage{siunitx}
\usepackage{siunitx}
\usepackage{siunitx}
\usepackage{amssymb}
\usepackage{tikz}
\usepackage{float}
\usepackage{siunitx}
\newtheorem{theorem}{Theorem}
\newtheorem{lemma}{Lemma}
 \def\INS#1#2#3#4 {\scalebox{0.8}{{\begin{tabular}{|c|c|c|c|} \hline $#1$ &$#2 $ & $#3$ & $#4$   \\ \hline \end{tabular}}}}
\newcommand{\R}{\mathbb{R}}
\newcommand{\md}{{\rm d}}
\title{Hadamard’s integral inequality as a tool for exposing latent convexity - Uniqueness results for polyconvex integrands}
\author{Jonathan Bevan, Martin Kru\v{z}\'ik, Jan Valdman}
\date{Source: arXiv:2604.09672v1 (CC BY 4.0)}
\begin{document}
\maketitle

\noindent\emph{Source and licence.} Hadamard’s integral inequality as a tool for exposing latent convexity - Uniqueness results for polyconvex integrands. Version arXiv:2604.09672v1 (CC BY 4.0), distributed by arXiv under the Creative Commons Attribution 4.0 International licence, http://creativecommons.org/licenses/by/4.0/, as recorded in the arXiv licence metadata for this exact version and shown on the abstract page https://arxiv.org/abs/2604.09672v1. Full licence notice: \texttt{licenses/arxiv-2604.09672-CC-BY-4.0.txt}.\par\smallskip

\noindent\textit{Editorial scope.} Counted unit: the article's lemma lowerbound (source0.tex lines 743--753). The article's complete original proof of it is delivered with it (source0.tex lines 755--785, one \texttt{proof} environment of 31 lines). No mathematics is written, repaired or re-derived: the statement and the proof are the article's own lines, in the article's own notation, and no retained line is substituted or re-flowed.  Retained uncounted as the source-local context the proof needs: lines 453--455.  Nothing was omitted from any retained span, so the package carries no omission record.  No retained line contains a citation, so the wrapper carries no bibliography and no bibliography entry.  No retained line contains a figure environment, an image-inclusion command or drawing code, so no image is delivered and the package declares no asset dependency.  Editorial formatting only, in addition: an AMS \texttt{amsart} preamble that restates the article's own declarations and macros used by the retained lines, namely the environments \texttt{theorem}, \texttt{lemma} on separate counters, together with the standard packages those lines need.  The retained environments print in this document as Theorem 1, Lemma 1 in order of appearance, whereas the article prints the same items with the numbering of its own full document.  The complete native source of the article as distributed in the arXiv e-print for this exact version is bundled unchanged as \texttt{sources/198221/source0.tex}.
\begin{theorem}\label{Thm:insulation}
Assume that $|c|\le 4$. Then $I\left(\varphi,\INS{-c}{0}{0}{c} \,\right)\ge 0$ for every $\varphi\in W^{1,2}_0(\Omega;\R^2)$, and the upper bound of $4$ is sharp.
\end{theorem}

\begin{lemma}\label{lowerbound}
 If $2<M\le 3$, $2/(M-2)\in \mathbb{N}$ and $\delta=(M-2)/4$
 then 
\begin{align}\label{special}
  \int_{R^\delta}|\nabla u_\varepsilon|^2 \,\md x  \ge \frac{M-2-\varepsilon}{2} \int_{R_{-2}}|\nabla u_\varepsilon|^2 \,\md x\ .
\end{align} 
 
for every $u_\varepsilon\in\mathcal{A}_\varepsilon$ defined above and 
every $0<\varepsilon<M-2$.
    
\end{lemma}

\begin{proof}
    
Assume that 
\begin{align}
  \int_{R^\delta}|\nabla u_\varepsilon|^2 \,\md x  < \frac{M-2-\varepsilon}{2} \int_{R_{-2}}|\nabla u_\varepsilon|^2 \,\md x\ .
\end{align}
This implies that 
\begin{align}
   \int_{R_{-2}}|\nabla u_\varepsilon|^2 \,\md x -M\int_{R_{-2}}\det\nabla u_\varepsilon\,\md x+ \int_{R^\delta}|\nabla u_\varepsilon|^2 \,\md x<0\ .
\end{align}

Let us further assume that $2/(M-2)\in\mathbb{N}$.
 We extend $u_\varepsilon$ from $R^\delta$ to $R_{-1}$ by 2/(M-2)  flips along the vertical direction   so that 
 \begin{align}\label{upper-bound}
     \int_{R_{-1}}|\nabla u_\varepsilon|^2 \,\md x=\frac{2}{M-2} \int_{R^\delta}|\nabla u_\varepsilon|^2 \,\md x<
     \frac{(M-2-\varepsilon)}{(M-2)}\int_{R_{-2}}|\nabla u_\varepsilon|^2 \,\md x
 \end{align}

 We have 
 \begin{align}\label{M=4}
    \int_{R_{-2}}|\nabla u_\varepsilon|^2 -4\int_{R_{-2}}\det\nabla u_\varepsilon\,\md x <  \int_{R_{-2}}|\nabla u_\varepsilon|^2 -\frac{4}{2+\varepsilon} \int_{R_{-2}}|\nabla u_\varepsilon|^2\, \md x=\frac{\varepsilon-2}{\varepsilon+2}\int_{R_{-2}}|\nabla u_\varepsilon|^2\, \md x\ .
 \end{align}

 Putting together \eqref{M=4} and \eqref{upper-bound}
we get for $2<M\le 3$
\begin{align}
  \int_{R_{-2}}|\nabla u_\varepsilon|^2 -4\int_{R_{-2}}\det\nabla u_\varepsilon\,\md x+ \int_{R_{-1}}|\nabla u_\varepsilon|^2 \,\md x< \frac{\varepsilon\,(2M-6-\varepsilon)}{(\varepsilon+2)(M-2)}\int_{R_{-2}}|\nabla u_\varepsilon|^2\, \md x <0\ ,
\end{align}
 which contradicts Theorem~\ref{Thm:insulation} whenever $2<M\le 3$. 

 \end{proof}

\end{document}
